% wfs_chapter_hda
% repository https://github.com/spatialaudio/wfs_chapter_hda
% authors:
% - Frank Schultz, https://orcid.org/0000-0002-3010-0294, https://github.com/fs446
% - Nara Hahn, https://orcid.org/0000-0003-3564-5864, https://github.com/narahahn
% - Sascha Spors, https://orcid.org/0000-0001-7225-9992, https://github.com/spors
%
% authors' versions for the chapters (English, German) on **Wave Field Synthesis** in
% - Stefan Weinzierl (editor): *Handbuch der Audiotechnik*, 2nd GER edition, Springer, 2025
% - Stefan Weinzierl (editor): *Handbook of Audio Technology*, 1st ENG edition, Springer, TBA
% - https://link.springer.com/book/10.1007/978-3-662-60369-7
%
% text and graphics under [CC BY 4.0 license](https://creativecommons.org/licenses/by/4.0/)
% source code under [MIT license](https://opensource.org/licenses/MIT)
% publisher Springer has copyright to their finally edited and authors' approved English / German chapters and their layouts
%
% we use the violine image from https://upload.wikimedia.org/wikipedia/commons/thumb/f/f1/Violin.svg/2048px-Violin.svg.png
% to create the files `python/violin_wfs_ENG.png` and `python/violin_wfs_DEU.png`
%
% we use the photo `fotos/WFS_Array_UniRostockH8_2014.jpg` CC BY 4.0 Matthias Geier & Sascha Spors
%
% all other graphics (as pdf, png, eps, svn, ipe) in this repository are CC BY 4.0 Frank Schultz & Nara Hahn
%
% we might also find https://git.iem.at/zotter/wfs-basics useful
%##############################################################################

%Wellenfeldsynthese (WFS) ist ein räumliches Wiedergabeverfahren, mit dem
%räumlich und zeitlich gestaffelte Wellenfronten mittels kontrollierter
%Interferenz synthetisiert werden.
Wave field synthesis (WFS) is a spatial reproduction technique.
It involves controlled interference to produce spatially and temporally
distinct wavefronts.
%
%Dazu benötigt es Lautsprecherarrays mit sehr dichtem Lautsprecherabstand
%und individueller Signalverarbeitung für die Lautsprecher.
This requires loudspeaker arrays with very dense loudspeaker spacing and
individual signal processing for each loudspeaker. 
%
%Im Gegensatz zu kanalbasierten Wiedergabeverfahren werden bei der WFS die
%Lautsprechersignale errechnet wofür Messdaten oder Audio-Objekte, und deren
%räumlich-zeitliche Parametrik einfließen.
Unlike channel-based reproduction methods, WFS calculates the loudspeaker
signals using measurement data or audio objects and their spatio-temporal
parameters.
%
%Anwendung findet die WFS bei lautsprecherbasierter Auralisation,
%Nachhallsynthese in Räumen, 3D-Beschallungskonzepten,
%audiologischer Grundlagenforschung und
%spatialisierter Audiokunst.
WFS is used in loudspeaker-based auralisation, reverberation enhancement,
3D sound reinforcement, audiological research and spatialised audio arts.

%Das Kapitel diskutiert die akustischen Grundlagen der WFS, enthält eine
%kompakte, anschauliche Herleitung der WFS einer virtuellen Punktquelle und
%arbeitet die wichtigsten Eigenschaften der mit WFS synthetisierten Schallfelder
%heraus.
%Abschließend werden praxisrelevante WFS-Modifikationen und -Werkzeuge dargelegt.
This chapter discusses the acoustic fundamentals of the WFS,
provides a compact and accessible derivation of the WFS for a
virtual point source, and highlights the key properties
of sound fields synthesised using the WFS.
Finally, it addresses practical WFS modifications and tools.